\documentclass[a4paper,12pt]{article}

% -----------------------------
% Pacchetti
% -----------------------------
\usepackage[ruled,vlined]{algorithm2e}
\usepackage{listings}
\usepackage{xcolor}
\usepackage{amsmath}
\usepackage{amssymb}

% -----------------------------
% Configurazione aspetto codice
% -----------------------------

\lstset{
	language=Python,
	basicstyle=\ttfamily\small,
	keywordstyle=\color{blue},
	stringstyle=\color{red},
	commentstyle=\color{green!50!black},
	numbers=left,
	numberstyle=\tiny\color{gray},
	frame=single,
	breaklines=true
}

% -----------------------------
% Definizione colori personalizzati
% -----------------------------
\definecolor{CommentBlue}{RGB}{30,90,200}
\definecolor{CommentGray}{RGB}{120,120,120}
\definecolor{CommentRed}{RGB}{200,40,40}

% -----------------------------
% Macro per commenti colorati
% -----------------------------
\newcommand{\BlueComment}[1]{\tcp{\textcolor{CommentBlue}{#1}}}
\newcommand{\GrayComment}[1]{\tcp{\textcolor{CommentGray}{#1}}}
\newcommand{\RedComment}[1]{\tcp{\textcolor{CommentRed}{#1}}}

% Versione a fine riga
\newcommand{\BlueCommentInline}[1]{\tcp*{\textcolor{CommentBlue}{#1}}}
\newcommand{\GrayCommentInline}[1]{\tcp*{\textcolor{CommentGray}{#1}}}
\newcommand{\RedCommentInline}[1]{\tcp*{\textcolor{CommentRed}{#1}}}



\title{Correspondence between NUTS implemented by hand and the Gellmann's article}
\author{Tommaso Pirazzo}
\date{}

\begin{document}
	
	\maketitle
	
	In the following there is the $6^{th}$ algorithm of the article \emph{"The No-U-Turn Sampler: Adaptively Setting Path Lengths
	in Hamiltonian Monte Carlo"} with the indicated correspondences with the code implemented in the file \verb|JAX_NUTS_functions.ipynb|.
	\vspace{.5cm}

\section{Leapfrog integrator}

	
\section{FindReasonableEpsilon}
	
	\begin{algorithm}[H]
		\caption{Heuristic for choosing an initial value of $\epsilon$}
		\KwIn{Current position $\theta$}
		\KwOut{Initial step size $\epsilon$}
		
		% Step 1: sample momentum
		Sample momentum $r \sim \mathcal{N}(0, I)$\;
		
		% Step 2: initial epsilon
		$\epsilon \leftarrow 1$\;
		
		% Step 3: compute initial joint log-density
		$E_0 \leftarrow \log p(\theta, r)$\;
		
		% Step 4: one leapfrog step
		$(\theta', r') \leftarrow \text{Leapfrog}(\theta, r, \epsilon)$\;
		$E_1 \leftarrow \log p(\theta', r')$\;
		
		% Step 5: acceptance statistic
		$\text{log\_accept} \leftarrow E_1 - E_0$\;
		
		% Step 6: choose direction
		\eIf{$\text{log\_accept} > \log(0.5)$}{
			$a \leftarrow +1$\;
		}{
			$a \leftarrow -1$\;
		}
		
		% Step 7: expand or shrink epsilon
		\While{$a \cdot \text{log\_accept} > -a \cdot \log(2)$}{
			$\epsilon \leftarrow \epsilon \cdot 2^{a}$\;
			$(\theta', r') \leftarrow \text{Leapfrog}(\theta, r, \epsilon)$\;
			$E_1 \leftarrow \log p(\theta', r')$\;
			$\text{log\_accept} \leftarrow E_1 - E_0$\;
		}
		
		\Return{$\epsilon$}\;
		
	\end{algorithm}

\begin{lstlisting}
def FindReasonableEpsilon(theta, key) :
 assert theta.ndim == 1
			
 # sample momentum r ~ N(0, I)
 key, subkey = jrnd.split(key)
 r = jrnd.normal(subkey, shape=theta.shape)
				
 # initial epsilon
 epsilon = 1.0
				
 # compute initial log joint density
 E0 = log_p(theta, r)
				
 # one leapfrog step
 theta_prime, r_prime = Leapfrog(theta, r, epsilon)
 E1 = log_p(theta_prime, r_prime)
				
 # log acceptance probability
 log_accept = E1 - E0
				
 # direction: +1 (increase) or -1 (decrease)
 a = jnp.where(log_accept > jnp.log(0.5), 1.0, -1.0)
				
 # -------------------------------
 # cond_fun: while loop condition
 # -------------------------------
 def cond_fun(state):
  eps, log_acc = state
  return (a * log_acc) > (-a * jnp.log(2.0))
				
 # --------------------
 # body_fun: loop body
 # --------------------
 def body_fun(state):
  eps, _ = state
  eps = eps * (2.0 ** a)
  theta_p, r_p = Leapfrog(theta, r, eps)
  E1 = log_p(theta_p, r_p)
  log_acc = E1 - E0
  return (eps, log_acc)
				
 # while loop in JAX
 epsilon, _ = jax.lax.while_loop(cond_fun, body_fun, (epsilon, log_accept))
				
 return epsilon
\end{lstlisting}

\section{BuildTree}
This pseudo-code is referred to the \verb*|BuildTree| function implemented in the \textbf{Algorithm 6} (No-U-Turn Sampler with Dual Averaging) of the article. This function is composed by two cases:
\begin{itemize}
	\item $j=0$: Base case $\rightarrow$ take one leapfrog step in the direction v.
	\item $j \neq 0$: Recursion $\rightarrow$ implicitly build the left and right subtrees.
\end{itemize}
The "base case" is implemented before in the function \verb*|build_tree_base| and then called in the complete function \verb*|BuildTree|. 

	\begin{algorithm}[H]
		\caption{BuildTree$(\theta, r, u, v, j, \epsilon)$}
		\KwIn{Position $\theta$, momentum $r$, slice variable $u$, direction $v \in \{-1,+1\}$, depth $j$, step size $\epsilon$}
		\KwOut{
			$(\theta^-, r^-)$, $(\theta^+, r^+)$, candidate $\theta'$, 
			number of valid points $n'$, stop flag $s'$, 
			acceptance sum $\alpha'$, acceptance count $n_\alpha'$
		}
		
		\If{$j = 0$}{
			% Base case: single leapfrog step
			$(\theta', r') \leftarrow \text{Leapfrog}(\theta, r, v \cdot \epsilon)$\;
			$E' \leftarrow \log p(\theta', r')$\;
			
			% Slice test
			$n' \leftarrow \mathbb{I}\left[u < \exp(E')\right]$\;
			
			% Stop criterion (turning)
			$s' \leftarrow \mathbb{I}\left[u < \exp(E')\right]$\;
			
			% Acceptance statistics
			$\Delta \leftarrow E' - \log p(\theta, r)$\;
			$\alpha' \leftarrow \min\{1, \exp(\Delta)\}$\;
			$n_\alpha' \leftarrow 1$\;
			
			\Return{
				$(\theta', r')$, $(\theta', r')$, $\theta'$, 
				$n'$, $s'$, $\alpha'$, $n_\alpha'$
			}\;
		}
		
		% Recursive case: build left subtree
		$(\theta^-, r^-), (\theta^+, r^+), \theta', 
		n', s', \alpha', n_\alpha' 
		\leftarrow \text{BuildTree}(\theta, r, u, v, j-1, \epsilon)$\;
		
		\If{$s' = 1$}{
			\Return{
				$(\theta^-, r^-), (\theta^+, r^+), \theta', 
				n', s', \alpha', n_\alpha'$
			}\;
		}
		
		% Build right subtree
		\eIf{$v = -1$}{
			$(\theta^-, r^-), (\_, \_), \theta'', 
			n'', s'', \alpha'', n_\alpha'' 
			\leftarrow \text{BuildTree}(\theta^-, r^-, u, v, j-1, \epsilon)$\;
		}{
			$(\_, \_), (\theta^+, r^+), \theta'', 
			n'', s'', \alpha'', n_\alpha'' 
			\leftarrow \text{BuildTree}(\theta^+, r^+, u, v, j-1, \epsilon)$\;
		}
		
		% Uniform choice of candidate
		$p \leftarrow \frac{n''}{n' + n''}$\;
		\If{\text{Bernoulli}(p)}{
			$\theta' \leftarrow \theta''$\;
		}
		
		% Combine statistics
		$n' \leftarrow n' + n''$\;
		$s' \leftarrow s' \cdot s'' \cdot 
		\text{Turning}(\theta^-, r^-, \theta^+, r^+)$\;
		$\alpha' \leftarrow \alpha' + \alpha''$\;
		$n_\alpha' \leftarrow n_\alpha' + n_\alpha''$\;
		
		\Return{
			$(\theta^-, r^-), (\theta^+, r^+), \theta', 
			n', s', \alpha', n_\alpha'$
		}\;
		
	\end{algorithm}
	

\section{Final algorithm}	
	\begin{algorithm}[H]
		\caption{No-U-Turn Sampler with Dual Averaging}
		Given $\theta_0$ , $\delta$, $L$, $M$, $M^{adapt}$ : \\
		\BlueComment{FindReasonableEpsilon($\theta$, key) $\rightarrow$ epsilon}
		Set $\epsilon_0$ = FindReasonableEpsilon($\theta$),\\ 
		 $\mu$ = $\log{10\epsilon_0}$, $\bar{\epsilon}_0$ = 1, $\bar{H}_0$ = 0, $\gamma$ = 0.05, $t_0$ = 10, $k$ = 0.75.
		 
		
		
		\KwIn{Array $A$ di lunghezza $n$, valore $x$}
		\KwOut{Indice $i$ tale che $A[i] = x$, oppure \texttt{None}}
		
		\For{$i \leftarrow 1$ \KwTo $n$}{
			\If{$A[i] = x$}{
				\Return $i$
			}
		}
		\Return None
	\end{algorithm}
	
\end{document}
FindReasonableEpsilon